\SourceLeaf{191}
\par\Needspace{6\baselineskip}\medskip
\begin{center}\hypertarget{native-191-article142}{}\pdfbookmark[2]{142.}{article-142}142.\end{center}
Equation IV, which would be of the eighth degree if expanded, is most readily solved by trial in its present form. The theory of equations also readily shows (we forgo a fuller development here for brevity) that this equation admits either two or four solutions in real values\HistoricalNoteEight{Manuscript note:] $z$ has at least two real values, since the values of $G\sin z^4-\sin(z+H)$ for $z=0$ and for $z=180^0$ have opposite signs. $z$ has no more than 4 real values, since $\dfrac{\sin(z+A)}{\sin z^4}$ has just one maximum and one minimum between $z=0$ and $z=180^0$, and likewise between $z=180^0$ and $360^0$. This occurs precisely when $2\operatorname{tang}z=-3\operatorname{cotang}A\pm\BareRoot{(9\operatorname{cotang}A^2-16)}=8\operatorname{tang}(z+A)$. —
\par\indent The solution of the equation $a\sin z^4=\sin(z-A)$ can be reduced to an elegant construction. Set $\operatorname{cotang}z=x$, so that $\sin z=\dfrac{1}{\BareRoot{(1+xx)}}\mathpunct{\cdot}$ Hence the equation
\[
\frac{1}{(1+xx)^{\frac{3}{2}}}=\frac{\cos A}{a}-\frac{\sin A}{a}x.
\]
The problem is thus to find the intersection of the curve $y=\dfrac{1}{(1+xx)^{\frac{3}{2}}}$ with the straight line $y=\dfrac{\cos A}{a}-\dfrac{\sin A}{a}x$.
\par\noindent\textsc{Binet}, \textit{Ecole polyt. Cah. 20. p. 285.}}. In the former case one value of $\sin z$ will be positive, and the negative one must be rejected, because by the nature of the problem $r'$ cannot be negative. In the latter case, either one value of $\sin z$ will be positive and the other three negative — so there can be no doubt which to adopt — or three will be positive and one negative. In this last case, any positive values for which $z$ is greater than $\delta'$ must also be rejected, because another essential condition of the problem requires $\rho'$, and therefore also $\sin(\delta'-z)$, to be positive.

When the observations are separated by moderate intervals of time, the last case will usually occur, so that three positive values of $\sin z$ satisfy the equation. Besides the true solution there is usually one in which $z$ differs little from $\delta'$, sometimes exceeding it and sometimes falling short. This phenomenon is explained as follows. The analytical treatment of our problem rests solely on the condition that the three positions of the celestial body in space must lie \textit{on} straight lines whose positions are determined by the absolute position of the Earth and the \SourceLeaf{192}observed position. Now the nature of the problem requires those positions to lie on \textit{those} parts of the straight lines \textit{from which} light reaches the Earth. The analytical equations, however, do not recognise this restriction and must include equally all sets of positions consistent with \textsc{Kepler}'s laws, whether they lie on these lines on one side of the Earth or the other, or coincide with the Earth itself. This last case will certainly satisfy our problem, since the Earth itself moves in accordance with those laws. It follows that the equations must include a solution in which the points $C,C',C''$ coincide with the points $A,A',A''$ (insofar as we neglect the very small changes that perturbations and parallax introduce into the Earth's elliptical positions). Equation IV would therefore always admit the solution $z=\delta'$ if the true values corresponding to the Earth's positions were taken for $P$ and $Q$. If, however, those quantities are assigned values that differ little from the true ones (as may always be assumed when the time intervals are short), there must necessarily be a solution of equation IV that comes close to $z=\delta'$.

Usually, when equation IV admits three solutions with positive values of $\sin z$, the third (besides the true one and the one just discussed) will give a value of $z$ greater than $\delta'$, and hence will be analytically possible but physically impossible. In that event there can be no doubt which to adopt. Yet the equation may admit two distinct admissible solutions, so that two entirely different orbits can satisfy our problem. In such a case the true orbit will easily be distinguished from the false one as soon as other observations farther removed in time can be examined.
